One annotator produced the study's labels. This appendix reports a check built
from a source no annotator touched, and the citation machinery behind
\S\ref{sec:long}.

\subsection{Judges label each other's opinions}
When a court writes ``\emph{Smith}, 123 F.3d 456, 460 (declining to decide
whether the statute applies),'' it asserts on the record that \emph{Smith} did
not resolve that question; ``(holding that \dots)'' asserts the opposite. We
scanned all \NEligible{} eligible opinions for reporter citations and applied a
frozen two-class characterization dictionary to a window around each.
\NCharLoose{} passages matched, but the loose form is unreliable: a marker near a
citation is often about something else, as where ``it did not reach a verdict''
refers to a jury. Requiring the verb to bind directly to the citation, as a
parenthetical or appositive, leaves \NCharTight{} passages, \NCharLZero{} of them
assertions of non-resolution, over \NCharOps{} cited opinions. We release this
set.

\subsection{Why it cannot validate the annotations}
The obvious use fails. Suppose a later court says \emph{Smith} held $X$ and our
annotation says \emph{Smith} left $X$ open. Two explanations fit: our label is
wrong, or the later court has overstated \emph{Smith}, which is the phenomenon
this paper measures. Nothing in the disagreement distinguishes them. We report
this because the design is tempting and the flaw shows only once the numbers are
in front of you.

\subsection{Proximity matching and why it failed}
Pairing each characterization with the sentence in the cited opinion sharing the
most content words gives \NJudgeBench{} items (\NJudgeUnres{} unresolved,
\NJudgeDec{} decided), but systems transfer poorly: the frozen dictionary reaches
\JudgeLexF{} macro-$F_1$ against 0.946 on the annotated benchmark. The cause is
measurable and is not the labels. Among items a judge called unresolved, only
8.2\% carry a dictionary expression in the selected anchor sentence, against
0.5\% of items a judge called decided: enriched sixteenfold, so the labels carry
signal, but in more than nine cases in ten the matcher selects where the opinion
\emph{discusses} the issue rather than where it \emph{disposes} of it. Accuracy
rises monotonically with match quality, from 0.485 to 0.610, the pattern
anchoring noise predicts and label noise does not. Asking instead whether a
dictionary expression sits within 400 characters of the characterized passage
covers \NEscN{} characterizations, of which \NProxyGood{} of \NProxyCheck{}
hand-checked flags were genuine. Both failures are the same failure, and it is
why \S\ref{sec:long} matches propositions.

\subsection{Chains, matching, and inference}
Absolute rates in \S\ref{sec:long} are not prevalence estimates, since the
characterization set is dominated by holding-ascriptions by construction, so only
contrasts are interpretable. The origin-side reading is high precision and low
recall, so the comparison runs between origins that \emph{visibly} mark an issue
open and the rest. Intervals resample originating opinions rather than passages,
since several citing passages share an origin. We fit no causal model: whether a
court leaves an issue open is a choice correlated with unobservables, including
how contested the question was, and those plausibly affect how later courts
describe it.

Locating where an opinion disposes of a named proposition, rather than where it
discusses one, is the binding constraint on a population rate. At the measured
yields, estimating one to within a few points needs on the order of a thousand
verified propositions, which at one in ten means annotating close to ten thousand
citing passages: an annotation budget an order of magnitude larger than a single
annotator supplies.
